The class \textbf{NP} is defined in the following way:
\begin{definition}
    A language $L \in \{0,1\}^*$ is in the class \textbf{NP} if and only if there exist a polynomial $p: N \rightarrow N$ and a polynomial time Turing Machine 
    $M$ such that:
    $$ L = \{x \in \{0,1\}^* | \exists y \in \{0,1\}^{p(|x|)}.M(\lfloor (x, y) \rfloor) = 1\} $$
    where:
    \begin{itemize}
        \renewcommand{\labelitemi}{-}
        \item $M \rightarrow$ \textbf{verifier} for language $L$
        \item $y \in \{0, 1\}^{p(|x|)} \text{ s.t. } M(\lfloor (x, y) \rfloor) = 1 \rightarrow$ \textbf{certificate} for $x$
    \end{itemize}
\end{definition}
\textit{NP} is the class of problems such that certificates are polynomial in length and they can be verified in polynomial time. \vspace{7pt}

Differently from \textit{P} and \textit{EXP} that have as counterpart the class of functions, the class \textit{NP} does not have a natural counterpart as 
a class of functions. 

Finally, we can state that \textit{NP} sits between \textit{P} and \textit{EXP}. More formally:
\begin{definition}[title = Theorem]
    $$ P \subseteq \textit{NP} \subseteq \textit{EXP} $$
\end{definition}
Every problem in \textbf{P} can be seen as a very special case of problem in \textbf{NP}, where the certificate $y$ is empty. Therefore, for every decision problem 
$L \in P$, there is a Turing Machine $N$ solving $L$ in polynomial time, without the need of any certificate. \vspace{7pt}

Formally, what we have discussed can be described as follows:
\begin{definition}
    If $M$ is a machine that on output $(\lfloor x, y \rfloor)$, accepts it if and only if $N$ accepts $y$, then:
    $$L = \{x \in \{0,1\}^*| \exists y \in \{0,1\}^{p(|x|)}.M(\lfloor x,y \rfloor) = 1 \}$$
    where:
    $p(|x|) = 0$. So, $y$ does not really exist.
\end{definition}